\theory{Model theory}{What is the relationship between a formal description and the mathematical structures in which it is true?}{Model theory studies languages, theories, structures, and definable sets. It asks which statements have models, how large those models can be, and how models compare.}{Mathematical logic, algebra, number theory, geometry, database theory, computability, and learning theory.}{Structures, formulas, and interpretations determine which properties are expressible and how algebraic models can be compared.}

\paragraph{The basic idea and history}
A formal language has symbols for functions, relations, and constants. A theory is a collection of sentences written in that language. A model is a structure in which every sentence of the theory is true. The integers, real numbers, groups, fields, and graphs can all be models of suitable theories \cite{model_theory_ref1}.

Alfred Tarski coined the term "theory of models" in 1954. Earlier work by Lowenheim, Skolem, Godel, and Malcev supplied the compactness and completeness tools. Since the 1970s, Saharon Shelah's stability theory has shaped the subject.

\end{historylist}
\section*{3. Theory}
\subsection*{Core theory}
\paragraph{Structures and formulas}
The language of rings has symbols $+$, $\times$, $-$, $0$, and $1$. In a field, the formula $y=x\times x$ defines the parabola $y=x^2$. A formula with free variables defines a set; a sentence with no free variables is either true or false in a structure.

Definable sets can use parameters from the model. Over the real numbers, $y=x^2+\pi$ is definable using $\pi$. Model theory studies not only the objects named by a theory but also which subsets can be described in its language.

\paragraph{Completeness, compactness, and size}
The completeness theorem says that a first-order sentence is true in every model of a theory exactly when it has a formal proof from that theory. The compactness theorem says that if every finite subset of a set of first-order sentences has a model, then the whole set has a model.

The downward Lowenheim--Skolem theorem says that a theory with an infinite model has models of many smaller infinite cardinalities, down to the size of its language. The upward theorem gives models in arbitrarily large cardinalities. Thus first-order axioms often cannot uniquely specify the size of an infinite structure.

\paragraph{Quantifier elimination}
Quantifiers such as "there exists" can hide the shape of a definable set. A theory has quantifier elimination if every formula is equivalent to one without quantifiers. Algebraically closed fields have quantifier elimination in the language of rings: every definable set is a Boolean combination of polynomial equations.

The theory of real closed fields also has quantifier elimination after adding an order relation. This yields the decidability of many elementary geometric questions and is the logical basis of cylindrical algebraic decomposition. Model completeness is a weaker property in which embeddings between models preserve all formulas in an appropriate sense \cite{model_theory_ref2}.

\paragraph{Ultraproducts and nonstandard analysis}
An ultraproduct combines a family of structures while ignoring a chosen collection of exceptional indices. Los's theorem says a first-order sentence holds in the ultraproduct exactly when it holds in almost all components according to the ultrafilter. Ultraproducts compare finite structures with infinite limits and help prove transfer principles.

Robinson used these ideas to build nonstandard analysis, where infinitesimal and infinitely large numbers are handled rigorously. Ax's work on pseudofinite fields and the Ax--Kochen theorem are important applications of ultraproduct methods.

\paragraph{Stability, o-minimality, and definable geometry}
Stability theory classifies theories by how complicated their definable sets and types can be. Stable, simple, NIP, and distal theories form part of a hierarchy. Geometric stability theory applies these ideas to algebraic and diophantine geometry.

O-minimal structures are ordered structures in which every definable subset of the line is a finite union of points and intervals. This tame one-dimensional behavior controls higher-dimensional definable geometry. Pila and others used o-minimal methods in problems related to the Andre--Oort conjecture. Hrushovski used model-theoretic geometry in the Mordell--Lang and Manin--Mumford settings \cite{model_theory_ref3}.

\paragraph{Applications}
Finite model theory studies finite structures and is central to database queries, descriptive complexity, and graph properties. A query is definable by a logical formula, and the expressive power of the logic can be compared with computational complexity classes.

Model theory also studies real algebraic geometry, differential fields, valued fields, p-adic fields, number fields, and pseudofinite fields. NIP theories and stability have connections to PAC learning and online learning classes. The theory is useful when algebraic structure and logical definability must be analyzed together.

\paragraph{What the theory depends on and where it is useful}
Model theory depends on first-order logic, set theory, algebra, and combinatorics. It helps algebra by transferring statements between models, geometry by describing definable sets, number theory by analyzing fields, and computer science by comparing logical descriptions with algorithms.

First-order logic cannot express every mathematical property, and a complete theory may have nonisomorphic models of different sizes. These are not failures; they identify exactly what the chosen language and axioms can control.

\section*{4. Important theorems and results}
\begin{enumerate}[leftmargin=*,itemsep=0.35em]
\item \textbf{Godel completeness theorem.} Semantic consequence in first-order logic equals formal provability.

\item \textbf{Compactness theorem.} If every finite part of a first-order theory is satisfiable, the entire theory is satisfiable.

\item \textbf{Lowenheim--Skolem theorems.} Infinite first-order theories have models in controlled smaller and larger cardinalities.

\item \textbf{Los's theorem.} First-order formulas in an ultraproduct are evaluated componentwise according to the ultrafilter.

\item \textbf{Tarski quantifier elimination.} Real closed fields admit quantifier elimination in the language of ordered rings \cite{model_theory_ref1}.
\end{enumerate}

\theoryapplications
\theoryconnections{model theory}{%
\item Logic supplies formulas and satisfaction, set theory supplies structures, and algebra supplies the groups, rings, and fields used as examples.
\item Compactness and ultraproducts connect syntax with global and asymptotic properties of structures.
}{%
\item Model theory helps algebra by transferring first-order statements between structures and helps computer science analyze formal languages and verification.
\item It also studies which properties can be expressed and decided, separating definability from computation.
}

\section*{Glossary}
\begin{description}[leftmargin=*,style=nextline,itemsep=0.3em]
\item[\textbf{Structure.}] A mathematical object consisting of a domain of elements together with named operations, relations, and constants that satisfy a given set of axioms.
\item[\textbf{Theory.}] A collection of sentences (statements with no free variables) in a formal language, all assumed to be true in any model of that theory.
\item[\textbf{Definable set.}] A subset of a structure that can be described by a first-order formula, possibly using fixed parameters from the structure.
\item[\textbf{Compactness theorem.}] If every finite subcollection of first-order sentences has a model, then the entire collection has a model, linking local consistency to global satisfiability.
\item[\textbf{Ultraproduct.}] A construction that builds a new structure from a family of structures by selecting which components to consider ``almost everywhere'' according to an ultrafilter.
\item[\textbf{Quantifier elimination.}] A property of a theory in which every formula is logically equivalent to one without existential or universal quantifiers, simplifying definable sets to Boolean combinations of basic equations.
\item[\textbf{First-order logic.}] A formal language with quantifiers over elements, connectives, and variables; the basic language of model theory.
\item[\textbf{Completeness theorem.}] A sentence is true in every model of a theory exactly when provable from that theory.
\item[\textbf{L\"owenheim--Skolem theorem.}] A first-order theory with an infinite model has models of every cardinality above the language size.
\item[\textbf{O-minimality.}] A tameness condition where every definable subset of the line is a finite union of points and intervals.
\item[\textbf{Stability theory.}] A classification of first-order theories by the complexity of their definable sets and types.
\end{description}

\section*{7. Sample Questions}
\emph{Exercise reference:} \cite{model_theory_ref1}. The cited edition is the source for the exercise set; consult its Exercises section for the official statement and numbering.
\begin{enumerate}[leftmargin=*,itemsep=0.9em]
\item \textbf{Exercise 1.} Let $\mathcal{L}=\{+\}$ and consider the structure $(\mathbb{Z},+)$. What is the smallest $\mathcal{L}$-substructure of $(\mathbb{Z},+)$ containing both $1$ and $-1$?

\emph{Solution.} A substructure must be closed under $+$ and contain $0$ (the identity). Starting from $\{1,-1\}$: $1+1=2$, $(-1)+(-1)=-2$, $1+(-1)=0$, and all integer multiples $n\cdot 1$ and $n\cdot(-1)$ are generated. The smallest such substructure is $(\mathbb{Z},+)$ itself.

\item \textbf{Exercise 2.} In the ultraproduct $\prod_{n\geq 2}\mathbb{Z}/n\mathbb{Z}\,/\,\mathcal{U}$ over a nonprincipal ultrafilter $\mathcal{U}$, how many solutions does the equation $x^2=x$ have?

\emph{Solution.} In each component $\mathbb{Z}/n\mathbb{Z}$, the equation $x^2=x$ has exactly $2$ solutions (when $n$ is prime: $x=0$ and $x=1$; for composite $n$ there may be more, but $0$ and $1$ always work). By Los's theorem, $x^2=x$ holds in the ultraproduct iff it holds for $\mathcal{U}$-almost all $n$. The sequences $\tilde{0}=(0,0,\ldots)$ and $\tilde{1}=(1,1,\ldots)$ are solutions. Any other solution $\tilde{x}$ must satisfy $x_n^2\equiv x_n\pmod{n}$ for $\mathcal{U}$-almost all $n$. Since $0$ and $1$ are the only idempotents in each $\mathbb{Z}/p\mathbb{Z}$ for prime $p$, and there are infinitely many primes in $\mathcal{U}$, we get exactly $2$ solutions: $\tilde{0}$ and $\tilde{1}$.

\item \textbf{Exercise 3.} In the structure $(\mathbb{Z},+,<,0,1)$, express the statement ``$n$ is even'' as a first-order formula with one free variable $x$.

\emph{Solution.} ``$x$ is even'' is $\exists y\,(x=y+y)$. More explicitly: $\exists y\,(x=y+y)$, which says there exists an integer $y$ with $x=2y$.

\item \textbf{Exercise 4.} Use the compactness theorem to prove that the theory $T=\mathrm{Th}(\mathbb{N},+,\cdot,0,1,<)$ (true sentences of arithmetic) has a nonstandard model---a model containing an element $c$ with $c\neq\bar{n}$ for every standard numeral $\bar{n}$.

\emph{Solution.} Add a new constant $c$ and consider $T'=T\cup\{c\neq\bar{n}:n\in\mathbb{N}\}$. Every finite subset of $T'$ involves finitely many sentences $c\neq\bar{n}_1,\ldots,c\neq\bar{n}_k$. The structure $(\mathbb{N},+,\cdot,0,1,<)$ interpreted with $c$ assigned to any integer larger than all $n_i$ satisfies this finite subset. By compactness, $T'$ has a model. This model satisfies $T$ (so it is a model of arithmetic) and contains an element $c$ different from every standard numeral.

\item \textbf{Exercise 5.} In the language of rings, write a sentence $\sigma_n$ expressing ``every polynomial of degree $n$ has a root'' over a field. Show that the theory of algebraically closed fields of characteristic $0$ is exactly $\{\sigma_n:n\geq 1\}\cup\{\text{char}=0\}$.

\emph{Solution.} $\sigma_n$ is $\forall a_0\cdots\forall a_n\,\bigl(a_n\neq 0\to\exists x\,(a_n x^n+\cdots+a_1 x+a_0=0)\bigr)$. A field is algebraically closed iff every polynomial of every degree has a root, so the theory of algebraically closed fields of characteristic $0$ is $\mathrm{ACF}_0=\{\sigma_n:n\geq 1\}\cup\{\neg\sigma_p^*:p\text{ prime}\}$ where $\sigma_p^*$ says $p\cdot 1=0$. This is exactly the set of sentences described.

\end{enumerate}
\section*{8. References}
\addcontentsline{toc}{section}{References}

\begin{thebibliography}{9}
\bibitem{model_theory_ref1} David Marker, \emph{Model Theory: An Introduction}, Springer, 2002.
\bibitem{model_theory_ref2} Wilfrid Hodges, \emph{A Shorter Model Theory}, Cambridge University Press, 1997.
\bibitem{model_theory_ref3} C. C. Chang and H. Jerome Keisler, \emph{Model Theory}, 3rd ed., North-Holland, 1990.
\end{thebibliography}